% ============================================================================
%  Part I -- The model and its prior
%  sections/part1_model_kernel.tex   (\input-ready fragment; no preamble)
%  Covers: (1) Introduction, (2) Generative model, (3) Arc-cosine kernel +
%  structured spatial prior.
% ============================================================================

\section{Introduction: the sequential active-learning problem}
\label{sec:intro}

\subsection{Setup}

We study a sequential active-learning problem for a Poisson--Gaussian-process
model of a latent function $\lat(x)$ over inputs $x$. An input
$x\in\Reals^{d}$ is a real-valued signal on a 2-D grid of $d$ grid points (e.g. a
$108\times108$ grid, so $d=108^2=11664$); grid point $i$ sits at grid coordinate
$\pixc_i$, and the input domain is a bounded box (e.g. $[0,1]^{d}$). For a single
input $x$ the model produces a scalar latent value $\lat(x)$, an exponential-link
rate $\rate(x)=\exp(\gain\,\lat(x)+\latz)$, and a Poisson count observation
$\spk\sim\Poi(\rate(x))$ (\S\ref{sec:genmodel}). None of the formulas below assume
the particular dimension $d=11664$.

The latent function $\lat$ is unknown a priori and must be inferred from the
counts. We place a Gaussian-process (GP) prior on $\lat$ and update its posterior
online as counts arrive. The setting is \emph{sequential}: after each observation
the model is refit and used to choose the next input, so the inputs presented
depend on the observations already collected.

\subsection{The active-learning goal}

The inputs are not presented in a fixed order. At each round the next input is
chosen to be maximally informative about $\lat$ --- this is \emph{active
learning}. A candidate input $x^{*}$ is scored by an information \emph{utility}
$\Ustd(x^{*})$ (Part~III) measuring the expected reduction in uncertainty about
$\lat$ if $x^{*}$ were observed next; the input maximizing the utility over the
available pool is chosen. A second variant $\Uda$ scores information about the rate
over the input (data) distribution $\px$. Because scoring and selecting from a
fixed pool is limited by the pool, a final stage (Part~IV) instead
\emph{synthesizes} an input that maximizes the utility, using a diffusion prior
over inputs steered by the utility gradient. All three stages consume two
derivatives of the GP prior kernel that are established in this part: the input
gradient $\nabla_x\Kf$~\eqref{eq:gradxK} (utility and diffusion guidance) and the
hyperparameter gradients $\partial\Cmat/\partial\theta$~\eqref{eq:dCtheta} (model
fitting).

\subsection{Roadmap of the document}

The document proceeds in the order the pipeline runs.

\begin{itemize}[leftmargin=1.4em]
\item \textbf{Part~I (this part) --- the model and its prior.} The generative model
  (\S\ref{sec:genmodel}): a latent GP over input space, an exponential link, and
  Poisson counts. The prior kernel (\S\ref{sec:kernel}): the order-1
  arc-cosine (ReLU) kernel built on a structured spatial prior covariance
  $\Cmat(\Amp,\bwid,\smoo,\rfc)$ that encodes a local, smooth weighting centred on
  the localized region, together with the two kernel gradients the rest of the
  pipeline needs.
\item \textbf{Part~II --- inference.} The Poisson likelihood breaks conjugacy, so
  the posterior over the latent is approximated by a sparse variational GP on $\Mind$
  inducing points with Gaussian posterior $q(\latt)=\Gauss(\vm,\vV)$, optimized
  against the evidence lower bound (ELBO); the model is trained by an EM-style
  coordinate-ascent loop (E-, F-, M-steps).
\item \textbf{Part~III --- active learning.} The information utility: the marginal
  count entropy $\Hmarg$, the production utility
  $\Ustd=\Hmarg-\E[\Hnoise]$, and the distribution-aware research variant
  $\Uda=\Hmarg-\E_{x\sim\px}[\Hcond]$; a deep layer of conditioning, subspace, and
  divergence results grounded in the kernel of this part.
\item \textbf{Part~IV --- synthesis.} Guided diffusion: a denoising diffusion prior
  whose reverse process is steered by $\nabla_x\Ustd$ (which flows through
  $\nabla_x\Kf$) to synthesize high-utility inputs.
\end{itemize}

\emph{The central open problem.} A single thread runs from the prior of this part
to the analysis of Part~III. The arc-cosine kernel is positively homogeneous in
the input: its self-value $\Kf(x,x)=x\transpose\Cmat x+\sigz^2$
(\S\ref{subsec:selfkernel}) grows with the input norm. Under utility maximization
the posterior variance then grows with the square of the input scale, so the
utility \emph{diverges} and the optimizer drives inputs to the boundary of the
bounded box (saturation). Whether a kernel exists that both (i) keeps the utility
bounded on the bounded input domain and (ii) fits well is the central \emph{open
problem}; it is set up here and analyzed in Part~III (with two candidate fixes, a
normalized arc-cosine kernel and a saturating arc-sin kernel).


\section{The generative model}
\label{sec:genmodel}

\subsection{Latent Gaussian process, exponential link, Poisson counts}

The count observation $\spk_i$ (also written $r_i$, or $n$ for a novel test input)
for input $x_i$ is a Poisson draw whose rate is an exponential link applied to a
latent Gaussian-process function $\lat$.

\begin{definition}[Generative model]
\label{eq:genmodel}
With a zero-mean GP prior over input space,
\begin{equation}
\lat(\cdot)\ \sim\ \GProc\!\left(0,\ \Kf\right),
\qquad
\rate(x)\ =\ \exp\!\big(\gain\,\lat(x)+\latz\big),
\qquad
\spk_i\ \sim\ \Poi\!\big(\rate(x_i)\big).
\end{equation}
Here $\lat(x)\in\Reals$ is the \emph{latent function} (what the GP learns;
\emph{not} the rate); $\rate(x)>0$ is the \emph{rate} (the Poisson rate
parameter); $\gain>0$ is the \emph{gain}, scaling how strongly the latent
modulates the rate; and $\latz\in\Reals$ is the \emph{bias} / log-baseline
rate, so that $\rate=e^{\latz}$ when $\lat=0$. The log-rate is
$\lgr(x)=\gain\,\lat(x)+\latz$, giving $\rate=e^{\lgr}$.
\end{definition}

For numerical stability the gain is constrained positive and bounded,
$\gain\in(0,10]$, and the bias to $\latz\in[-50,50]$.

The GP prior is defined by the property that any finite set of latent values
$\{\lat(x_1),\dots,\lat(x_n)\}$ is jointly Gaussian, $\lat(X)\sim\Gauss(0,\Kf)$
with $[\Kf]_{ij}=\Kf(x_i,x_j)=\mathrm{Cov}[\lat(x_i),\lat(x_j)]$. The mean is
zero. The kernel $\Kf$ is fixed in \emph{form} before seeing data but its entries
depend on the inputs through evaluation; its full definition is
\S\ref{sec:kernel}.

\subsection{Why the GP models the latent, not the rate}

Modelling $\lat$ as Gaussian and pushing it through $\rate=\exp(\gain\lat+\latz)$
means the rate is \emph{log-normally} distributed: we learn a scaled, shifted log
of the Poisson rate rather than the rate itself. Two properties make this the
right choice for count data.

\begin{itemize}[leftmargin=1.4em]
\item \emph{Positivity for free.} The exponential link guarantees $\rate>0$ for any
  real latent value. A Gaussian prior placed \emph{directly} on the rate would
  assign mass to negative (invalid) rates.
\item \emph{Smoothness where it belongs.} The GP prior imposes spatial structure
  (locality and smoothness, \S\ref{sec:kernel}) on the unconstrained latent
  $\lat$, not on the constrained rate. The prior belief --- that $\lat$ is a
  spatially local, smooth function of the input --- is naturally expressed on the
  unconstrained log-rate, which is where the structure lives.
\end{itemize}

\begin{remark}[Poisson breaks conjugacy]
For a Gaussian likelihood the posterior and predictive would be closed-form. With
the Poisson likelihood the marginal
$p(\spk\mid X,\theta)=\int \Poi(\spk\mid\rate(x))\,\Gauss(\lat\mid0,\Kf_\theta)\,d\lat$
has no closed form, so neither the exact posterior $p(\lat\mid\spk,X,\theta)$ nor
the exact predictive can be computed directly. This is what forces the variational
treatment of Part~II; here we fix only the model and its prior.
\end{remark}


\section{The arc-cosine kernel and the structured spatial prior}
\label{sec:kernel}

The prior covariance $\Kf$ is a first-order arc-cosine (ReLU) kernel built on a
structured input-covariance matrix $\Cmat$. This section owns the definitions of
both $\Kf$ and $\Cmat$, their parameterization and constraints, and the two
gradients ($\nabla_x\Kf$, $\partial\Cmat/\partial\theta$) that the utility, the
diffusion guidance, and the M-step consume.

\subsection{The order-1 arc-cosine kernel}
\label{subsec:acos}

\begin{definition}[Order-1 arc-cosine kernel]
For two input vectors $x,x'\in\Reals^{d}$,
\begin{equation}
\Kf(x,x')\ =\ \frac{1}{\pi}\,\Kmag\,\Jang(\kang),
\qquad
\Kmag=\sqrt{\vmag_x\,\vmag_{x'}},
\qquad
\cos\kang=\frac{x\transpose\Cmat x'+\sigz^2}{\Kmag},
\qquad
\Jang(\kang)=\sin\kang+(\pi-\kang)\cos\kang,
\label{eq:acoskernel}
\end{equation}
with self-magnitudes (prior variances) $\vmag_x=x\transpose\Cmat x+\sigz^2$ and
$\vmag_{x'}={x'}\transpose\Cmat x'+\sigz^2$, cross-term
$x\transpose\Cmat x'+\sigz^2$, magnitude $\Kmag$, angle $\kang$, and order-1
angular term $\Jang$. Equivalently, expanding the angle,
\begin{equation}
\kang=\arccos\!\left(
\frac{x\transpose\Cmat x'+\sigz^2}
     {\sqrt{x\transpose\Cmat x+\sigz^2}\ \sqrt{{x'}\transpose\Cmat x'+\sigz^2}}
\right).
\end{equation}
Here $\sigz^2$ is a constant \emph{bias variance} and $\Cmat\in\Reals^{d\times d}$
is the structured spatial prior covariance (\S\ref{subsec:Cprior}).
\end{definition}

\begin{remark}[Order and ReLU-network interpretation]
This is the order-$n=1$ member of the Cho--Saul (2009) arc-cosine family
($n=0$ is Heaviside/step, $n=1$ is ReLU); the angular term
$\Jang(\kang)=\sin\kang+(\pi-\kang)\cos\kang$ is exactly the ReLU case. The kernel
is the covariance of the output of an infinite-width, single-hidden-layer ReLU
network whose input-to-hidden weights are drawn $w\sim\Gauss(0,\Cmat)$, with the
bias contributing $\sigz^2$:
\begin{equation}
\Kf(x,x')\ =\ \E_{w\sim\Gauss(0,\Cmat)}\!\big[\,\mathrm{ReLU}(w\transpose x)\,\mathrm{ReLU}(w\transpose x')\,\big]
\quad\text{(up to the $1/\pi$ constant, with bias $\sigz^2$).}
\end{equation}
Thus $\Cmat$ is literally the prior covariance of the network weights across grid
points. Shaping $\Cmat$ (\S\ref{subsec:Cprior}) injects the prior belief that
$\lat$ depends on the input through a spatially \emph{local} and \emph{smooth}
region of grid points.
\end{remark}

\subsection{Non-stationarity and the self-kernel identity}
\label{subsec:selfkernel}

\begin{proposition}[Self-kernel identity]
The prior variance at any input is the $\Cmat$-quadratic form plus the bias
variance:
\begin{equation}
\Kf(x,x)\ =\ \vmag_x\ =\ x\transpose\Cmat x+\sigz^2.
\label{eq:selfkernel}
\end{equation}
\end{proposition}

\begin{proof}
At $x'=x$ the cross-term equals $\vmag_x=\Kmag$, so $\cos\kang=1$, $\kang=0$, and
$\Jang(0)=\sin0+(\pi-0)\cos0=\pi$. Hence
$\Kf(x,x)=\Kmag\,\Jang(0)/\pi=\Kmag=\sqrt{\vmag_x\,\vmag_x}=\vmag_x$.
\end{proof}

\begin{remark}[Non-stationarity and its consequence for active learning]
Equation~\eqref{eq:selfkernel} shows $\Kf$ depends on the actual input values
through $x\transpose\Cmat x$, not only on $x-x'$: the kernel is \emph{non-stationary}
and the prior variance $\vmag_x$ grows with the input norm (in the $\Cmat$-metric).
This has a direct downstream effect: the active-learning utility is driven by
posterior variance, so utility maximization is biased toward high-norm inputs
unless the kernel is normalized. For exactly this reason there are normalized and
saturating siblings of the base kernel --- a normalized arc-cosine kernel
($\bar{\Kf}=\Jang(\kang)/\pi$, constant $\bar{\Kf}(x,x)=1$) and a saturating
arc-sin kernel (bounded diagonal) --- analyzed in Part~III as fixes to the
norm-scaling divergence. This part keeps the focus on the base arc-cosine kernel.
\end{remark}

\subsection{The structured spatial prior covariance $\Cmat$}
\label{subsec:Cprior}

$\Cmat$ encodes that the localized support of the prior is spatially local and
smooth. It is \emph{constant with respect to the input} $x$ --- it depends only on
the hyperparameters $(\Amp,\bwid,\smoo,\rfc)$. Grid point $i$ sits at grid
coordinate $\pixc_i$ on a normalized $[-1,1]\times[-1,1]$ grid, and
$\rfc=(\varepsilon_{0x},\varepsilon_{0y})$ is the center of the localized region.

\begin{definition}[Structured spatial prior]
With an amplitude $\Amp$, a per-grid-point locality weight $\loc_i$, and a pairwise
smoothing factor $\Csm$,
\begin{equation}
\Cmat_{ij}
\ =\ \Amp\,\loc_i\,[\Csm]_{ij}\,\loc_j
\ =\ \Amp\,\exp\!\left(-\frac{\lVert\pixc_i-\rfc\rVert^2}{4\bwid^2}\right)\,
      \exp\!\left(-\frac{\lVert\pixc_i-\pixc_j\rVert^2}{2\smoo^2}\right)\,
      \exp\!\left(-\frac{\lVert\pixc_j-\rfc\rVert^2}{4\bwid^2}\right),
\label{eq:Cprior}
\end{equation}
where
\begin{equation}
\begin{aligned}
\loc_i&=\exp\!\left(-\frac{\lVert\pixc_i-\rfc\rVert^2}{4\bwid^2}\right)
&&\text{(locality weight)},\\
[\Csm]_{ij}&=\exp\!\left(-\frac{\lVert\pixc_i-\pixc_j\rVert^2}{2\smoo^2}\right)
&&\text{(smoothness kernel)},
\end{aligned}
\end{equation}
followed by a symmetrization $\Cmat\leftarrow(\Cmat+\Cmat\transpose)/2$.
\end{definition}

\emph{Interpretation of the natural parameters} (defaults $\Amp=1.0$,
$\bwid=\smoo=0.1$ on the $[-1,1]$ grid):
\begin{itemize}[leftmargin=1.4em]
\item $\Amp$ --- overall \emph{amplitude}: a positive gain (default $1.0$) on the
  structured covariance $\Cmat$, scaling the prior variance of the latent. It
  multiplies the $\Cmat$-quadratic part of the self-magnitude
  $\vmag_x=x\transpose\Cmat x+\sigz^2$ (not the bias $\sigz^2$); because it enters
  the kernel inside $\Kmag=\sqrt{\vmag_x\vmag_{x'}}$ and $\cos\kang$, it rescales
  $\Kf$ non-linearly rather than as a plain output scale. Unlike $\bwid,\smoo,\rfc$
  it sets the magnitude, not the spatial shape, of the localized weighting.
\item $\bwid$ --- \emph{half-width} of the localized region: the locality weight
  $\loc$ is a Gaussian of standard deviation $\sqrt{2}\,\bwid$ in grid-distance,
  reaching $e^{-1}$ at $\text{dist}=2\bwid$; for $\bwid=0.1$ the localized region
  spans roughly $\pm0.2$ ($\approx20\%$ of the axis).
\item $\smoo$ --- smoothness \emph{standard deviation} (correlation length) of the
  RBF-like $\Csm$ (variance $\smoo^2$): grid points closer than $\smoo$ are
  strongly correlated, grid points beyond $3\smoo$ nearly uncorrelated.
\item $\rfc$ --- the 2-D center of the localized region on $[-1,1]^2$.
\end{itemize}
Together, $\loc$ localizes weight around $\rfc$ and $\Csm$ correlates nearby grid
points --- a localized, smooth weight prior.

\subsection{Raw parameterization and constraints}
\label{subsec:rawparam}

For optimizer stability the model is parameterized by unconstrained log-space
\emph{raw} parameters rather than $\bwid$ and $\smoo$ directly:
\begin{equation}
\begin{aligned}
\text{raw\_m2log2beta}&=-2\log(2\bwid)&&\Rightarrow&&\exp(\text{raw\_m2log2beta})=\frac{1}{4\bwid^2},\\
\text{raw\_mlog2rho2}&=-\log(2\smoo^2)&&\Rightarrow&&\exp(\text{raw\_mlog2rho2})=\frac{1}{2\smoo^2},
\end{aligned}
\end{equation}
and $\sigz=\exp(\text{raw\_sigma\_0})$ under a positivity constraint (so the bias
variance added in~\eqref{eq:acoskernel} is $\sigz^2=(\exp(\text{raw\_sigma\_0}))^2$).
The exponentials of the raw parameters are used \emph{directly} as the exponent
coefficients, so that
$\loc_i=\exp(-\tfrac{1}{4\bwid^2}\lVert\pixc_i-\rfc\rVert^2)$ and
$[\Csm]_{ij}=\exp(-\tfrac{1}{2\smoo^2}\lVert\pixc_i-\pixc_j\rVert^2)$ reproduce
exactly the natural-parameter form~\eqref{eq:Cprior} (numerically,
$\bwid=\smoo=0.1\Rightarrow 1/(4\bwid^2)=25$, $1/(2\smoo^2)=50$). The raw
parameters invert back to the natural values,
$\bwid=\tfrac12\exp(-\text{raw\_m2log2beta}/2)$,
$\smoo=\tfrac{1}{\sqrt2}\exp(-\text{raw\_mlog2rho2}/2)$, and round-trip cleanly.

\begin{table}[h]
\centering
\begin{tabular}{@{}llll@{}}
\toprule
Param & Constraint mechanism & Bounds (natural) & Bounds (raw) \\
\midrule
$\sigz$ & $\sigz=\exp(\text{raw\_sigma\_0})$, positive & $\sigz>0$ & --- \\
$\bwid$ & via $\text{raw\_m2log2beta}$, clamped & $[0.01,\,0.3]$ & $[1.02,\,7.82]$ \\
$\smoo$ & via $\text{raw\_mlog2rho2}$, clamped & $[0.01,\,0.5]$ & $[0.69,\,8.52]$ \\
$\rfc$ & direct clamp & $[-1,1]$ per axis & --- \\
$\Amp$ & $\Amp=\mathrm{softplus}(\text{raw\_Amp})$, clamp & $(0,\,1000]$ & --- \\
\bottomrule
\end{tabular}
\end{table}

\begin{sloppypar}
Bounds are enforced two ways: by read-only rejection of an out-of-bounds trial
step during optimization, and by a projected clamp applied after each optimizer
step.
\end{sloppypar}

\begin{remark}[Restricting $\Cmat$ to its effective support]
The locality weight $\loc_i=\exp(-\lVert\pixc_i-\rfc\rVert^2/(4\bwid^2))$ decays as
a Gaussian away from $\rfc$, so grid points whose weight $\loc_i$ falls below a
small threshold (e.g. $10^{-3}$) contribute negligibly to $x\transpose\Cmat x$ and
may be dropped, shrinking $\Cmat$ from $(d,d)$ to $(d',d')$ with $d'\ll d$ (the
same restriction is applied to the input so only the retained grid points enter
$x\transpose\Cmat x$). With threshold $10^{-3}$ the retained support is the disk
$\lVert\pixc_i-\rfc\rVert\le2\bwid\sqrt{\ln1000}\approx5.25\,\bwid$ around $\rfc$;
the retained entries of $\Cmat$ still carry gradients, while the support itself is
computed from detached parameters and is therefore held fixed
(non-differentiable) and stable.
\end{remark}

\subsection{Kernel gradient with respect to the input, $\nabla_x\Kf$}
\label{subsec:gradxK}

Choosing the next input (and steering diffusion, Part~IV) requires the derivative
of the kernel with respect to the input. Treating $\Cmat$ as symmetric and
constant, $\sigz^2$ constant, and $x'$ (hence $\vmag_{x'}$) constant, the useful
intermediate gradients are
\begin{equation}
\nabla_x\,\vmag_x=2\,\Cmat x,
\qquad
\nabla_x\,(x\transpose\Cmat x'+\sigz^2)=\Cmat x',
\qquad
\nabla_x\,\Kmag=\sqrt{\tfrac{\vmag_{x'}}{\vmag_x}}\,\Cmat x,
\end{equation}
and the angular derivative $d\Jang/d\kang=-(\pi-\kang)\sin\kang$.

\begin{proposition}[Input gradient of the arc-cosine kernel]
\begin{equation}
\nabla_x\Kf(x,x')
\ =\ \frac{1}{\pi}\left[(\pi-\kang)\,\Cmat x'
\ +\ \sin\kang\,\sqrt{\frac{\vmag_{x'}}{\vmag_x}}\,\Cmat x\right]
\ =\ \frac{1}{\pi}\left[(\pi-\kang)\,\Cmat x'
\ +\ \sin\kang\,\sqrt{\frac{{x'}\transpose\Cmat x'+\sigz^2}{x\transpose\Cmat x+\sigz^2}}\,\Cmat x\right].
\label{eq:gradxK}
\end{equation}
By symmetry ($x\leftrightarrow x'$),
\begin{equation}
\nabla_{x'}\Kf(x,x')
=\frac{1}{\pi}\left[(\pi-\kang)\,\Cmat x
+\sin\kang\,\sqrt{\frac{x\transpose\Cmat x+\sigz^2}{{x'}\transpose\Cmat x'+\sigz^2}}\,\Cmat x'\right].
\end{equation}
\end{proposition}

\begin{proof}
Write $\Kf=\tfrac1\pi\Kmag\,\Jang(\kang)$ and apply the product rule
$\nabla_x\Kf\propto(\nabla_x\Kmag)\Jang+\Kmag(\nabla_x\Jang)$. Using
$\nabla_x\Jang=(d\Jang/d\kang)\nabla_x\kang$ with
$\nabla_x\kang=-(1/\sin\kang)\,\nabla_x(\cos\kang)$ and
$d\Jang/d\kang=-(\pi-\kang)\sin\kang$, the $\sin\kang$ cancels and
$\nabla_x\Jang=(\pi-\kang)\,\nabla_x(\cos\kang)$ with
$\nabla_x(\cos\kang)=\tfrac1\Kmag\Cmat x'-\tfrac{\cos\kang}{\Kmag}\sqrt{\vmag_{x'}/\vmag_x}\,\Cmat x$
(recall $\cos\kang=(x\transpose\Cmat x'+\sigz^2)/\Kmag$). Collecting terms,
\[
\pi\,\nabla_x\Kf
=\underbrace{\sin\kang\sqrt{\tfrac{\vmag_{x'}}{\vmag_x}}\Cmat x}_{[A]}
+\underbrace{(\pi-\kang)\cos\kang\sqrt{\tfrac{\vmag_{x'}}{\vmag_x}}\Cmat x}_{[B]}
+\underbrace{(\pi-\kang)\Cmat x'}_{[C]}
-\underbrace{(\pi-\kang)\cos\kang\sqrt{\tfrac{\vmag_{x'}}{\vmag_x}}\Cmat x}_{[D]} .
\]
Terms $[B]$ and $[D]$ are identical with opposite sign and cancel, leaving
$[A]+[C]$, which is~\eqref{eq:gradxK}.
\end{proof}

Every term carries a factor $\Cmat$, so $\nabla_x\Kf$ lives in the column space of
$\Cmat$; this is the algebraic root of the subspace analysis in Part~III (gradient
ascent on the utility already concentrates on the top $\Cmat$-eigenvectors).

\subsection{Kernel gradients with respect to the hyperparameters, $\partial\Cmat/\partial\theta$}
\label{subsec:dCtheta}

The M-step (Part~II training) maximizes the ELBO over the kernel hyperparameters and
needs $\partial\Cmat/\partial\theta$, which then chains through
$\vmag_x=x\transpose\Cmat x$ and the cross-term into $\partial\Kf/\partial\theta$
using the same quadratic-form structure as~\S\ref{subsec:gradxK} with $\Cmat$
replaced by $\partial\Cmat/\partial\theta$.

\begin{proposition}[Structured-prior hyperparameter gradients (natural parameters)]
Writing $\Cmat=\Amp\,\loc\,\Csm\,\loc\transpose$, the amplitude enters linearly and
the shape parameters through their Gaussian exponents. Each shape-parameter
derivative is $\Cmat_{ij}$ --- which already carries the $\Amp$ factor --- times the
derivative of its log-argument, while the amplitude derivative strips one factor of
$\Amp$:
\begin{equation}
\label{eq:dCtheta}
\begin{aligned}
\frac{\partial\Cmat_{ij}}{\partial\Amp}
&=\loc_i\,[\Csm]_{ij}\,\loc_j=\frac{\Cmat_{ij}}{\Amp},\\[2pt]
\frac{\partial\Cmat_{ij}}{\partial\bwid}
&=\Cmat_{ij}\,\frac{\lVert\pixc_i-\rfc\rVert^2+\lVert\pixc_j-\rfc\rVert^2}{2\bwid^3},\\[2pt]
\frac{\partial\Cmat_{ij}}{\partial\smoo}
&=\Cmat_{ij}\,\frac{\lVert\pixc_i-\pixc_j\rVert^2}{\smoo^3},\\[2pt]
\nabla_{\rfc}\Cmat_{ij}
&=\frac{\Cmat_{ij}}{2\bwid^2}\big[(\pixc_i-\rfc)+(\pixc_j-\rfc)\big]\qquad(\text{2D vector}).
\end{aligned}
\end{equation}
\end{proposition}

\begin{proof}
For $\Amp$: $\Cmat=\Amp\,(\loc\,\Csm\,\loc\transpose)$ is linear in $\Amp$, so
$\partial\Cmat_{ij}/\partial\Amp=\loc_i[\Csm]_{ij}\loc_j=\Cmat_{ij}/\Amp$. For $\bwid$:
$\partial_\bwid\big[-(\lVert\pixc_i-\rfc\rVert^2+\lVert\pixc_j-\rfc\rVert^2)/(4\bwid^2)\big]
=+(\lVert\pixc_i-\rfc\rVert^2+\lVert\pixc_j-\rfc\rVert^2)/(2\bwid^3)$. For $\smoo$:
$\partial_\smoo\big[-\lVert\pixc_i-\pixc_j\rVert^2/(2\smoo^2)\big]=+\lVert\pixc_i-\pixc_j\rVert^2/\smoo^3$.
For $\rfc$: $\nabla_{\rfc}\lVert\pixc-\rfc\rVert^2=2(\rfc-\pixc)$, giving
$\nabla_{\rfc}\log\Cmat_{ij}=(1/(2\bwid^2))[(\pixc_i-\rfc)+(\pixc_j-\rfc)]$.
\end{proof}

\begin{remark}[Gradients in the raw parameterization]
Because the model optimizes the raw parameters (\S\ref{subsec:rawparam}), the
corresponding \emph{raw}-parameter gradients follow from~\eqref{eq:dCtheta} by the
chain rule:
\begin{equation}
\begin{aligned}
\frac{\partial\Cmat_{ij}}{\partial(\text{raw\_m2log2beta})}
&=\Cmat_{ij}\,(\log\loc_i+\log\loc_j)
=-\Cmat_{ij}\,\frac{\lVert\pixc_i-\rfc\rVert^2+\lVert\pixc_j-\rfc\rVert^2}{4\bwid^2},\\
\frac{\partial\Cmat_{ij}}{\partial(\text{raw\_mlog2rho2})}
&=\Cmat_{ij}\,\log[\Csm]_{ij}
=-\Cmat_{ij}\,\frac{\lVert\pixc_i-\pixc_j\rVert^2}{2\smoo^2}.
\end{aligned}
\end{equation}
and, since $\rfc$ is a direct parameter,
$\partial\Cmat_{ij}/\partial\rfc=(\Cmat_{ij}/(2\bwid^2))[(\pixc_i-\rfc)+(\pixc_j-\rfc)]$
coincides with the natural form in~\eqref{eq:dCtheta}. The raw and natural
gradients are related by the constraint-transform Jacobian
$d(\text{raw})/d\bwid=-2/\bwid$, etc.
\end{remark}

\subsection{Numerical stability}
\label{subsec:numstab}

Three stability mechanisms act at three distinct levels; they must not be conflated.

\begin{enumerate}[leftmargin=1.6em]
\item \textbf{$\Cmat$ matrix --- symmetrization only.}
  $\Cmat\leftarrow(\Cmat+\Cmat\transpose)/2$ cancels the floating-point rounding
  asymmetry from the outer product. \emph{No diagonal jitter is added to $\Cmat$.}
\item \textbf{Kernel evaluation --- domain clamps.} With
  $\text{eps}=10^{-7}$, $\cos\kang$ is clamped into $(-1,1)$ so $\arccos$ is finite,
  and $1-\cos^2\kang$ is clamped at $\min=\text{eps}$ so $\sin\kang=\sqrt{1-\cos^2\kang}$
  never takes the square root of a negative rounding result. These keep the
  transcendental operations in-domain and are unrelated to Cholesky jitter.
\item \textbf{GP covariance / Cholesky --- jitter on the gram matrices, not on
  $\Cmat$.} Jitter is added to the \emph{output} gram matrices of the arc-cosine
  kernel --- the inducing gram $\Ktil=\Kf(\Ipt,\Ipt)$ and the predictive gram
  $\Kf(X,X)$ --- \emph{not} to $\Cmat$. A pre-Cholesky jitter (default $10^{-4}$) is
  added to $\Ktil$ and $\Kf(X,X)$; the kernel evaluation itself adds none (that
  would double-jitter). On a failed factorization the gram matrix is promoted to
  double precision (sequential Cholesky subtractions are prone to catastrophic loss
  of accuracy in single precision) and the factorization is retried with escalating
  jitter $10^{-4},10^{-3},10^{-2}$ (up to three tries).
\end{enumerate}

\caveat{Some intuition expects diagonal jitter on the prior $\Cmat$; there is none.
$\Cmat$ receives only symmetrization. Jitter lives exclusively on the GP gram
matrices $\Ktil$/$\Kf(X,X)$. Do not ``add jitter to $\Cmat$'' thinking it is
missing.}
